\subsection{A conical full subcategory opposite a small extreme corner}

\begin{proposition}[Proved]\label{CP2:A-conical-block-cut}
Let $C$ be a finite category with an object $p$ such that no
nonidentity arrow between distinct objects enters $p$, or no
such arrow leaves $p$. Suppose that $|\operatorname{End}_C(p)|\leq4$
and that the full deletion $E=C\setminus\{p\}$ is nonempty
and has an initial or a terminal object. Then every product
of positive-degree rational cohomology classes of $BC$ is zero.
In particular $BC$ cannot have the homotopy type of $\mathbb{CP}^2$.

For a category with acyclic support and hom-cardinality matrix
$h_{ij}$, it suffices that an extreme object $p$ has
$h_{pp}\leq4$ and that there exists $q\ne p$ with either
\[
 h_{qj}=1\quad\hbox{for every }j\ne p,
 \qquad\hbox{or}\qquad
 h_{jq}=1\quad\hbox{for every }j\ne p.
\]
These conditions include $h_{qq}=1$.
\end{proposition}
\begin{proof}
Let $D$ be the full one-object subcategory on $p$. The
small-monoid theorem, Corollary~\ref{CP2:A-five-types}, makes
$BD$ rationally acyclic. An actual initial object of $E$
contracts its nerve: the unique arrows from it to every
object define a natural transformation from its constant
functor to the identity. A terminal object gives the dual
contraction. Thus $BE$ is contractible.

If no strict arrow enters $p$, the partition $D,E$ has no
arrows from $E$ to $D$. Apply Lemma~\ref{CP2:A-two-block} to
this actual partition. If no strict arrow leaves $p$, apply
the same lemma with the blocks reversed. In either case,
every positive-degree rational class restricts to zero on
both blocks. The double-mapping-cylinder model has relative
quotient $\Sigma(T_+)$ after collapsing its two endpoint
spaces. Each class therefore lifts to relative cohomology,
where all positive-degree products are zero. Its absolute
product is consequently zero as well. This proves the claim
without a bounded-homology assumption or a condition on the
connecting profunctor beyond the specified direction.

For the cardinality criterion, $h_{qj}=1$ for all objects
$j$ of $E$ is precisely the assertion that there is a unique
arrow $q\to j$, including the identity when $j=q$. It makes
$q$ an actual initial object independently of every remaining
composition entry and of the other corner monoids. The column
condition similarly gives an actual terminal object. Thus
the displayed test proves the obstruction for every category
with the indicated profile, without enumerating its compositions.
\end{proof}

\begin{remark*}
The absent arrows between $p$ and some objects of $E$ cause no
problem: a connecting profunctor may have empty components.
The proof requires only the one-way support cut. It does not
delete $p$ or claim that its deletion inclusion is a homotopy
equivalence. For example, in the frozen four-object eighteen-arrow
rows $1773$ and $1774$, take $p=1$ and $q=0$. The deletion on
$0,2,3$ has initial object $0$ because $h_{00}=h_{02}=h_{03}=1$,
while $p$ is a trivial source. Their six-element corner at $2$
is unrestricted by this argument.
\end{remark*}
